\documentclass[11pt]{article}

\usepackage[T1]{fontenc}
\usepackage{lmodern}
\usepackage{amsmath,amssymb,amsthm}
\usepackage[margin=1in]{geometry}
\usepackage{booktabs}
\usepackage{enumitem}
\usepackage{microtype}
\usepackage[hidelinks]{hyperref}

\hypersetup{
  pdftitle={Far-left stability of a chemotaxis front: a mathematical proof of the sharp threshold chi=4}
}

\newcommand{\R}{\mathbb{R}}
\newcommand{\norm}[1]{\left\lVert #1\right\rVert}
\newcommand{\abs}[1]{\left\lvert #1\right\rvert}
\setlength{\emergencystretch}{3em}

\newtheorem{theorem}{Theorem}
\newtheorem{proposition}{Proposition}
\newtheorem{lemma}{Lemma}
\newtheorem{remark}{Remark}

\title{Far-Left Stability of a Chemotaxis Front\\
\large A mathematical proof of the sharp threshold \(\chi=4\), with the
necessary basin condition}
\author{}
\date{27 July 2026}

\begin{document}
\maketitle

\begin{abstract}
For the normalized one-dimensional chemotaxis--logistic equation
\[
 u_t=u_{zz}+(c-\chi v_z)u_z
      +u\bigl(1-\chi v+(\chi-1)u\bigr),
 \qquad -v_{zz}+v=u,
\]
we identify the exact stability threshold of the plateau
\((u,v)=(1,1)\).  The linear dispersion relation is strictly stable exactly
for \(0\le\chi<4\).  A nonlinear logarithmic entropy has the same sharp
coercive coefficient \(1-\chi^2/16\).  A slowly varying integrable weight
turns this identity into a whole-line Liouville theorem, and a
translation-compactness argument gives uniform far-left convergence for
orbits that remain in a persistent positive far-left band.

The band hypothesis is necessary for a broad class of front-like solutions.
A Fisher--KPP traveling front observed from a faster frame remains positive
at the left end of every fixed time slice, yet converges to zero at each
fixed point of that frame.  Thus slice-by-slice positivity cannot replace a
time-uniform basin or phase condition.
\end{abstract}

\tableofcontents

\section{Equation and far-left convergence}

In the normalized case \(m=\alpha=\gamma=1\), write the equation in a frame
moving at speed \(c\):
\begin{align}
 u_t
 &=u_{zz}+cu_z-\chi(uv_z)_z+u(1-u),                    \label{eq:div}\\
 -v_{zz}+v&=u.                                         \label{eq:ell}
\end{align}
Putting \(r=v-1\), equation \eqref{eq:ell} becomes
\[
 -r_{zz}+r=u-1.
\]
Since \(v_{zz}=v-u\), expanding the chemotactic divergence in
\eqref{eq:div} gives the equivalent equation
\[
 u_t=u_{zz}+(c-\chi v_z)u_z
      +u\bigl(1-\chi v+(\chi-1)u\bigr).
\]

A traveling front has the end states
\[
 (u,v)=(1,1)\quad\text{as }z\to-\infty,
 \qquad
 (u,v)=(0,0)\quad\text{as }z\to+\infty.
\]
The desired far-left conclusion for a time-dependent orbit is
\begin{equation}
 \forall\varepsilon>0\ \exists R,T\
 \forall t\ge T\ \forall z\le-R,\qquad
 \abs{u(t,z)-1}<\varepsilon.                            \label{eq:goal}
\end{equation}

\section{The exact linear threshold}

Set \(u=1+p\), \(v=1+r\).  The linearized equations are
\[
 p_t=p_{zz}+cp_z+(\chi-1)p-\chi r,
 \qquad -r_{zz}+r=p.
\]
For \(p=e^{\lambda t+ikz}\) and \(s=k^2\), the elliptic equation gives
\(\widehat r=\widehat p/(1+s)\), hence
\begin{equation}
 \lambda(s)=-s+ick+(\chi-1)-\frac{\chi}{1+s}.           \label{eq:lambda}
\end{equation}
The real growth rate is
\begin{equation}
 d_\chi(s)=\operatorname{Re}\lambda(s)
 =-1-s+\frac{\chi s}{1+s}.                             \label{eq:disp}
\end{equation}

\begin{theorem}[sharp dispersion threshold]
For \(0\le\chi<4\), \(d_\chi(s)<0\) for every \(s\ge0\).  At \(\chi=4\),
the modes \(k=\pm1\) are neutral.  For \(\chi>4\), a nonempty band of
Fourier modes grows exponentially.
\end{theorem}

\begin{proof}
If \(0\le\chi\le1\), then
\[
 d_\chi'(s)=-1+\frac{\chi}{(1+s)^2}\le0,
\]
so the maximum is \(d_\chi(0)=-1\).  If \(\chi>1\), the unique critical
point is \(s_*=\sqrt\chi-1\), and \(d_\chi''(s)<0\), so it is the global
maximum.  Substitution gives
\[
 d_\chi(s_*)=\chi-2\sqrt\chi.
\]
This is negative, zero, or positive according as \(\chi<4\), \(\chi=4\),
or \(\chi>4\).  At equality \(s_*=1\), hence \(k=\pm1\).
\end{proof}

More generally, a reaction exponent \(\alpha>0\) and coupling
\(\chi\gamma\ge0\) give
\[
 d_{\alpha,\chi\gamma}(s)
 =-\alpha-s+\frac{\chi\gamma\,s}{1+s}.
\]
When \(\chi\gamma>1\), its maximum is
\((\sqrt{\chi\gamma}-1)^2-\alpha\).  The exact stability condition is
\[
 \chi\gamma<(1+\sqrt\alpha)^2.
\]

\section{The sharp nonlinear entropy}

For \(u>0\), define
\[
 h(u)=u-1-\log u,\qquad h'(u)=1-\frac1u,
\]
and abbreviate
\[
 A=\frac{u_z}{u},\qquad W=u-1.
\]

\begin{lemma}[unweighted entropy identity]
On a periodic interval, or on the whole line when all boundary terms vanish,
\begin{equation}
 \frac{d}{dt}\int h(u)
 =-\int A^2+\chi\int A r_z-\int W^2.                   \label{eq:raw}
\end{equation}
\end{lemma}

\begin{proof}
Multiplying \eqref{eq:div} by \(h'(u)\) gives
\[
 \int h'(u)u_{zz}=-\int h''(u)u_z^2=-\int A^2.
\]
The constant drift is the integral of a derivative.  For the chemotactic
term,
\begin{align*}
 -\chi\int h'(u)(uv_z)_z
 &= -\chi\int\left(1-\frac1u\right)
       (u_zr_z+ur_{zz})\\
 &= -\chi\int\bigl(W_zr_z+Wr_{zz}-Ar_z\bigr)
  =\chi\int Ar_z.
\end{align*}
The first two terms cancel after integration by parts.  Finally,
\[
 \int h'(u)u(1-u)=\int (u-1)(1-u)=-\int W^2.
\]
\end{proof}

\begin{lemma}[sharp elliptic resolver estimate]
If \(-r_{zz}+r=W\), then
\begin{equation}
 \int r_z^2\le\frac14\int W^2.                         \label{eq:quarter}
\end{equation}
The constant \(1/4\) is optimal.
\end{lemma}

\begin{proof}
Taking Fourier transforms gives
\[
 \widehat{r_z}(k)=\frac{ik}{1+k^2}\widehat W(k).
\]
By Plancherel,
\[
 \int r_z^2
 =\int\frac{k^2}{(1+k^2)^2}\abs{\widehat W(k)}^2\,dk.
\]
The multiplier is at most \(1/4\), with equality when \(k^2=1\).
\end{proof}

Young's inequality gives
\[
 \chi Ar_z\le A^2+\frac{\chi^2}{4}r_z^2.
\]
Combining this with \eqref{eq:raw} and \eqref{eq:quarter} yields
\begin{equation}
 \frac{d}{dt}\int h(u)
 \le-\left(1-\frac{\chi^2}{16}\right)\int(u-1)^2.
                                                               \label{eq:sharp}
\end{equation}
No smallness assumption on \(u-1\) was used.  The nonlinear entropy margin
is positive exactly for \(\abs\chi<4\), matching the linear threshold.

\section{Slow localization on the whole line}

The unweighted entropy need not be integrable for a bounded entire solution.
For \(0<k<1\) and a center \(x_0\), introduce
\begin{equation}
 \phi_{k,x_0}(x)
 =\exp\!\left[-k\sqrt{1+(x-x_0)^2}\right].             \label{eq:weight}
\end{equation}
Then \(\phi>0\), \(\phi\in L^1(\R)\), and
\[
 \abs{\phi_x}\le k\phi.
\]

\subsection{A near-sharp weighted resolver bound}

Let
\[
\begin{aligned}
 S&=\int\phi W^2,&
 A_0&=\int\phi r^2,&
 I&=\int\phi r_z^2,\\
 C_0&=\int\phi r_{zz}^2,&
 J&=\int\phi_xrr_z,&
 P&=\int\phi Wr.
\end{aligned}
\]
Testing \(-r_{zz}+r=W\) by \(\phi r\) and expanding the square of the
equation give
\begin{align}
 A_0+I+J&=P,                                           \label{eq:test}\\
 2P&\le S+A_0,                                         \label{eq:pair}\\
 S&=A_0+2I+2J+C_0.                                    \label{eq:square}
\end{align}
Moreover,
\[
 \abs J\le\frac{k}{2}(A_0+I),
\]
and the nonnegativity of
\(\int\phi(r+r_{zz})^2\) gives
\[
 2I+2J\le A_0+C_0.
\]

\begin{lemma}[weighted resolver]
For \(0\le k<1\),
\begin{equation}
 I\le G(k)S,\qquad
 G(k)=\frac{1+k}{4(1-k)}.                              \label{eq:gain}
\end{equation}
\end{lemma}

\begin{proof}
Equations \eqref{eq:test}--\eqref{eq:pair} and the lower bound on \(J\)
give
\[
 (1-k)(A_0+I)\le S.
\]
Equation \eqref{eq:square} and \(2I+2J\le A_0+C_0\) give
\[
 4I+4J\le S.
\]
Hence
\[
 4I\le S+4\abs J
 \le S+2k(A_0+I).
\]
Multiplying by \(1-k\) and using
\((1-k)(A_0+I)\le S\) yields
\[
 4(1-k)I\le(1+k)S,
\]
which is \eqref{eq:gain}.
\end{proof}

The gain \(G(k)\) tends to the sharp value \(1/4\) as \(k\downarrow0\).

\subsection{The explicit localization loss}

Assume that \(u\ge\ell>0\).  The elementary inequality
\begin{equation}
 h(u)\le\frac{(u-1)^2}{u}\le\frac{(u-1)^2}{\ell}        \label{eq:hbound}
\end{equation}
follows from \(1-u^{-1}\le\log u\).  Weighted integration by parts, reserving
a fraction \(k\) of the diffusion to control the weight flux, gives
\begin{equation}
 \frac{d}{dt}\int_\R\phi h(u)
 \le-\bigl(1-\mathcal L(\chi,c,\ell,k)\bigr)
       \int_\R\phi(u-1)^2,                             \label{eq:localized}
\end{equation}
where
\begin{equation}
\begin{aligned}
 \mathcal L(\chi,c,\ell,k)
 ={}&\frac{\chi^2G(k)}{4(1-k)}
     +\frac{k}{4}
     +\frac{\chi k}{2}\bigl(1+G(k)\bigr)
     +\frac{\abs c\,k}{\ell}.                          \label{eq:loss}
\end{aligned}
\end{equation}

For completeness, the two estimates containing all weight losses are
\begin{align*}
 \chi X
 &\le(1-k)D+\frac{\chi^2}{4(1-k)}I,\\
 F
 &\le kD+
 \left(\frac{k}{4}+\frac{\chi k}{2}
              +\frac{\abs c\,k}{\ell}\right)S
 +\frac{\chi k}{2}I.
\end{align*}
Here
\[
 D=\int\phi\left(\frac{u_z}{u}\right)^2,\quad
 X=\int\phi\frac{u_z}{u}r_z,\quad
 S=\int\phi(u-1)^2,
\]
and \(F\) is the sum of terms in which a derivative falls on \(\phi\).
The first inequality is Young's inequality.  The second follows from
\(\abs{\phi_x}\le k\phi\), \eqref{eq:hbound}, and Young's inequality on
each diffusion and chemotactic flux.  Inserting \(I\le G(k)S\) gives
\eqref{eq:localized}--\eqref{eq:loss}.

At \(k=0\),
\[
 \mathcal L(\chi,c,\ell,0)=\frac{\chi^2}{16}.
\]
The loss is continuous at \(k=0\).  Therefore, for every
\(0\le\chi<4\), every \(c\in\R\), and every \(\ell>0\), one can choose
\begin{equation}
 0<k<1,\qquad \mathcal L(\chi,c,\ell,k)<1.              \label{eq:kchoice}
\end{equation}

\section{A whole-line Liouville theorem}

\begin{theorem}[entropy rigidity]
Let \((u,r)\) be a bounded classical solution on all
\((t,x)\in\R^2\) of
\[
 u_t=u_{xx}+cu_x-\chi(ur_x)_x+u(1-u),
 \qquad -r_{xx}+r=u-1.
\]
Assume \(0\le\chi<4\) and
\[
 0<\ell\le u(t,x)\le M
 \qquad\text{for all }(t,x)\in\R^2,
\]
together with the regularity and boundedness needed for the weighted
integrations above.  Then \(u\equiv1\) and \(r\equiv0\).
\end{theorem}

\begin{proof}
Choose \(k\) as in \eqref{eq:kchoice}, fix \(x_0\), and define
\[
 E(t)=\int_\R\phi_{k,x_0}(x)h(u(t,x))\,dx.
\]
The bounded range and \(\phi\in L^1\) give a uniform bound
\[
 0\le E(t)\le K\qquad(t\in\R).
\]
By \eqref{eq:hbound} and \eqref{eq:localized},
\[
 E'(t)\le-\rho E(t),\qquad
 \rho=\bigl(1-\mathcal L(\chi,c,\ell,k)\bigr)\ell>0.
\]
Starting at \(t-T\) and applying Gronwall's inequality gives
\[
 0\le E(t)\le E(t-T)e^{-\rho T}\le Ke^{-\rho T}.
\]
Letting \(T\to\infty\) yields \(E(t)=0\).  Since \(\phi>0\),
\(h(u)\ge0\), and \(h(u)=0\) only at \(u=1\), continuity implies
\(u(t,x)=1\) for every \(t,x\).  The elliptic equation then gives the
bounded solution \(r\equiv0\).
\end{proof}

The use of negative time is legitimate because the theorem concerns an
entire translated limit.  This is the decisive advantage of combining
entropy with late-time translation compactness.

\section{Far-left convergence by compactness}

We state the remaining analytic input without any hidden notation.
Assume:
\begin{enumerate}[label=\textup{(H\arabic*)},leftmargin=2.8em]
\item there are \(T,Z,\ell,M\), with \(\ell>0\), such that
\begin{equation}
 \ell\le u(t,z)\le M
 \qquad(t\ge T,\ z\le Z);                              \label{eq:band}
\end{equation}
\item for every pair of sequences \(t_n\to+\infty\) and
  \(z_n\to-\infty\), the translates
\[
 u_n(t,x)=u(t+t_n,x+z_n),\qquad
 r_n(t,x)=r(t+t_n,x+z_n)
\]
and the fields
\[
 (u_n)_x,\ (u_n)_{xx},\ (u_n)_t,\ r_n,\ (r_n)_x,\
 (r_n)_{xx}
\]
are locally equicontinuous.  Moreover, there is a constant \(C\), independent
of the chosen sequences and of the compact box, such that all seven fields
have absolute value at most \(C\) on each fixed compact box for all
sufficiently large \(n\).  Consequently a common subsequence converges
locally uniformly to globally bounded fields, with the differential
identities passing to the limit.
\end{enumerate}

\begin{theorem}[basin-conditional far-left convergence]
If \(0\le\chi<4\) and \emph{(H1)--(H2)} hold, then
\eqref{eq:goal} holds.
\end{theorem}

\begin{proof}
Suppose \eqref{eq:goal} fails.  Then there are
\(\varepsilon_0>0\), \(t_n\to\infty\), and \(z_n\to-\infty\) such that
\[
 \abs{u(t_n,z_n)-1}\ge\varepsilon_0.
\]
By \emph{(H2)}, a subsequence of the translates converges locally uniformly
to a bounded entire classical solution \((U,R)\).  Hypothesis
\eqref{eq:band} passes to the limit and gives
\(\ell\le U\le M\) on all of \(\R^2\).  The differential equations also
pass to the limit.  The Liouville theorem therefore gives \(U\equiv1\).
On the other hand, local convergence at the origin gives
\[
 \abs{U(0,0)-1}
 =\lim_n\abs{u(t_n,z_n)-1}
 \ge\varepsilon_0,
\]
a contradiction.
\end{proof}

\section{Why the basin condition is necessary}

Slice-by-slice left positivity has the quantifier pattern
\[
 \forall t\ \exists\ell_t>0\ \exists Z_t\
 \forall z\le Z_t,\qquad u(t,z)\ge\ell_t.
\]
The persistent band \eqref{eq:band} instead requires
\[
 \exists\ell>0\ \exists T,Z\
 \forall t\ge T\ \forall z\le Z,\qquad u(t,z)\ge\ell.
\]
The order of the quantifiers cannot be interchanged.

\begin{proposition}[convective escape]
At \(\chi=0\), there is a genuine global classical front which is positive
at the left end of every fixed time slice but has no persistent positive
floor when observed from a sufficiently fast frame.
\end{proposition}

\begin{proof}
For every \(s\ge2\), the Fisher--KPP equation has a monotone traveling front
\(U_s\) satisfying
\[
 U_s''+sU_s'+U_s(1-U_s)=0,\qquad
 U_s(-\infty)=1,\quad U_s(+\infty)=0.
\]
Take \(s=3\), and let \(V_s\) be the bounded solution of
\(-V_s''+V_s=U_s\).  In a frame of speed \(c=4\), set
\[
 q(t,z)=U_s\bigl(z+(c-s)t\bigr),\qquad
 w(t,z)=V_s\bigl(z+(c-s)t\bigr).
\]
The chain rule and the wave equation give
\[
 q_t=q_{zz}+cq_z+q(1-q),\qquad -w_{zz}+w=q,
\]
which is precisely the \(\chi=0\) moving-frame system.  For each fixed
\(t\), \(q(t,z)\to1\) as \(z\to-\infty\).  But for each fixed \(z\),
\[
 z+(c-s)t\to+\infty,\qquad q(t,z)\to0.
\]
Given any proposed \(\ell>0,T,Z\), choose \(t\ge T\) so large that
\(q(t,Z)<\ell\).  Thus no persistent band exists.
\end{proof}

This counterexample does not contradict stability relative to a fixed wave
under a weighted closeness or phase condition: such a condition pins the
solution to the same wave speed.  It rules out only an unconditional
far-left theorem over arbitrarily phased front-like orbits.

\section{Conclusion}

The number \(4\) is sharp in three compatible senses:
\begin{enumerate}[leftmargin=2em]
\item below \(4\), every Fourier mode decays; at \(4\), \(k^2=1\) is
  neutral; above \(4\), a band of modes grows;
\item the nonlinear entropy has the exact coefficient
  \(1-\chi^2/16\);
\item throughout \(0\le\chi<4\), slow localization, the entire-orbit
  Liouville theorem, and translation compactness yield uniform far-left
  convergence once a persistent positive band is supplied.
\end{enumerate}
The band or an equivalent wave-phase condition is structural, not technical:
without it, a genuine front can leave a faster observation frame by
convective escape.

\end{document}
